\input{../article_base.tex} \title{פתרון מטלה 09 --- משוואות דיפרנציאליות, 80320}
% chktex-file 9
% chktex-file 17

\DeclareMathOperator\Div{div}

\usepackage{pgfplots, tikz, pgffor}
\usetikzlibrary{math}
\pgfplotsset{compat=1.18}

\begin{document}
\maketitle
\maketitleprint[blue]{}

\question{}
תהי מערכת שטורם־ליוביל,
\[
	\frac{d}{d t} (p(t) \frac{d}{d t} u(t)) + (q(t) + \lambda) u(t) = 0,
	\quad
	u(a) = u(b) = 0
.\]
עבור $p > 0$.

\subquestion{}
יהי $V = C([a, b])$ מרחב הפונקציות הרציפות ונגדיר מכפלה פנימית על $V$ על־ידי,
\[
	\langle u, v \rangle
	= \int_{a}^{b} u(x) v(x)\ dx
.\]
ונגדיר את $W = C_0^1([a, b])$ מרחב הפונקציות הגזירות ברציפות ב־$[a, b]$ המתאפסות בקצוות,
ונגדיר את $D : W \to V$ על־ידי $D u = u'$. \\
נראה ש־$D$ אופרטור אנטי־סימטרי, כלומר,
\[
	\forall u, v \in W,
	\quad
	\langle D u, v \rangle
	= - \langle u, D v \rangle
.\]
\begin{proof}
	תהיינה $f, g \in W$, אז מתקיים,
	\[
		\langle D f, g \rangle
		= \int_{a}^{b} f'(x) g(x)\ dx
		= \left. f(x) g(x) \right\rvert_a^b - \int_{a}^{b} f(x) g'(x)\ dx
		= - \langle f, D g \rangle
	.\]
	על־ידי משפט אינטגרציה בחלקים,
	כאשר $\left. f(x) g(x) \right\rvert_a^b = 0$ שכן $f(a) = f(b) = 0$ מההנחה כי $f \in W$.
\end{proof}

\subquestion{}
הפעם נסמן $V = C([a, b]), W = C_0^2([a, b])$ מרחב הפונקציות הגזירות ברציפות פעמיים המתאפסות בקצוות, ויהי $\Delta = D^2$ אופרטור הגזירה פעמיים.
נראה ש־$\Delta : W \to V$ הוא אופרטור סימטרי ונסיק כי האופרטור $S : W \to V$ המוגדר על־ידי,
\[
	S u(t)
	= \frac{d}{d t} (p(t) \frac{d}{d t}u(t)) + q(t) u(t)
.\]
הוא אופרטור סימטרי גם כן.
\begin{proof}
	תהיינה $f, g \in C_0^2([a, b])$ ושוב נחשב,
	\[
		\langle \Delta f, g \rangle
		= \langle D^2 f, g \rangle
		= -\langle D f, D g \rangle
		= {(-1)}^2 \langle f, D^2 g \rangle
		= \langle f, \Delta g \rangle
	.\]
	כאשר השתמשנו אנטי־סימטריה שמצאנו בסעיף הקודם והעובדה שאם $f \in C_0^2([a, b])$ אז בפרט $f \in C_0^1([a, b])$.

	אנו מניחים כי $p, q$ גזירות ברציפות פעמיים, וכן $C_0^2([a, b])$ סגורה לכפל בפונקציה גזירה ברציפות פעמיים, לכן כהרכבת פונקציות ושני האופרטורים מהסעיפים האחרונים נסיק שגם $S$ סימטרית.
\end{proof}

\subquestion{}
נוכיח שאם $T : W \to V$ אופרטור סימטרי ו־$u, v \in W$ וקטורים עצמיים עם ערך עצמי שונה אז $u \perp v$. \\
נסיק שאם $u, v$ פותרות את המערכת עבור ערכים עצמיים שונים אז,
\[
	\int_{a}^{b} u(t) v(t)\ dt
	= 0
.\]
\begin{proof}
	נסמן $T u = \lambda u, T v = \mu v$ ולכן $\lambda \ne \mu$.
	מסימטריה,
	\[
		\lambda \langle u, v \rangle
		= \langle T u, v \rangle
		= \langle u, T v \rangle
		= \mu \langle u, v \rangle
	.\]
	נקבל אם כך ש־$(\lambda + \mu) \langle u, v \rangle = 0$ ובפרט מההנחה שמדובר בווקטורים עצמיים של ערכים עצמיים שונים נוכל להניח ש־$\lambda + \mu \ne 0$ (אחרת נקבל עד כדי היפוך הווקטור ערך עצמי זהה) ש־$u \perp v$.

	בסעיף הקודם מצאנו ש־$S$ סימטרי, אז אם $u, v$ פתרונות עבור $\lambda, \mu$ שונים אז בפרט המכפלה הפנימית שהגדרנו בסעיף א' מתאפסת וקיבלנו את המבוקש.
\end{proof}

\question{}
תהיינה $p, q$ פונקציות גזירות ברציפות כך ש־$p > 0$ ונגדיר את הבעיה,
\[
	\frac{d}{d t} (p(t) \frac{d}{d t} u(t)) + (q(t) + \lambda) u(t) = 0,
	\quad
	u(a) = u(b) = 0
.\]
לכל $\lambda \in \RR$ נגדיר את $u_{\lambda}(t)$ להיות הפתרון לבעיה,
\[
	\frac{d}{d t} (p(t) \frac{d}{d t} u(t)) + (q(t) + \lambda) u(t) = 0,
	\quad
	u(a) = 0,
	u'(a) = 1
.\]
וניזכר בהגדרת הקורדינטות הפולריות,
\[
	p(t) u_{\lambda}'(t) = r_{\lambda}(t) \cos \theta_{\lambda}(t),
	\quad
	u_{\lambda}(t) = r_{\lambda}(t) \sin \theta_{\lambda}(t)
.\]
נראה ש־$\lambda \mapsto \theta_{\lambda}(b)$ היא פונקציה רציפה.
\begin{proof}
	נזכור שמצאנו בכיתה ש־$\theta_{\lambda}$ מקיימת את המשוואה הבאה,
	\[
		\theta_{\lambda}'(t)
		= (q(t) + \lambda) \sin^2 \theta_{\lambda}(t) + \frac{1}{p(t)} \cos^2 \theta_{\lambda}(t),
	.\]
	כלומר זוהי פונקציה שמקיימת משוואה דיפרנציאלית הומוגנית וגזירה ברציפות, ולכן אם נגדיר,
	\[
		F(x, y, \lambda)
		= (q(y) + \lambda) \sin^2(x) + \frac{1}{p(t)} \cos^2(x)
	.\]
	אז $\theta_{\lambda}'(t) = F(\theta_{\lambda}(t), t, \lambda)$ ובפרט קיימת $U$ פתוחה מקסימלית כך ש־$\Phi(\tau_{\lambda}^0, t)$ זרימה, ולכן בפרט ממשפט רציפות זרימות היא רציפה בשלושת משתניה.
	נעיר כי את המשפט הוכחנו עבור משוואות הומוגניות, אך משקילות נקבל נכונות גם במקרה זה, ובפרט נקבל רציפות במשתנה השלישי בהצבה במשתנה השני $b$, כלומר $\lambda \mapsto \theta_{\lambda}(b)$ רציפה.
\end{proof}

\question{}
בשאלה זו נוכיח שבקורדינטות פולריות מתקיימת הנוסחה,
\[
	\Div F
	= \frac{1}{r} \partial_r (r F_r) + \frac{1}{r} \partial_{\varphi} F_{\varphi}
.\]
יהי $F : \RR^2 \to \RR^2$ שדה וקטורי גזיר, נסמן $\hat{x} = (1, 0), \hat{y} = (0, 1)$ וקטורי היחידה, ונסמן את הפירוק של $F$,
\[
	F = F_x \hat{x}  + F_y \hat{y}
.\]
נגדיר קורדינטות גליליות על־ידי,
\[
	x = r \cos \varphi,
	\quad
	y = r \sin \varphi
.\]
ונגדיר את וקטורי היחידה של המערכת הגלילית על־ידי $\hat{r}, \hat{\varphi}$ ונבחין כי הם מקיימים,
\[
	\hat{r}(x, y)
	= \frac{(x, y)}{\lVert (x, y) \rVert}
	= \frac{x \hat{x} + y \hat{y}}{\sqrt{x^2 + y^2}},
	\qquad
	\hat{\varphi} = \frac{-y \hat{x} + x \hat{y}}{\sqrt{x^2 + y^2}}
.\]
וכן נסמן גם,
\[
	F = F_r \hat{r} + F_{\varphi} \hat{\varphi}
.\]
נסמן גם $\partial_x, \partial_y$ נגזרות במערכת הקרטזית וכן $\partial_r, \partial_{\varphi}$ נגזרות במערכת הגלילית.

\subquestion{}
נראה כי,
\[
	F_r
	= \cos \varphi F_x + \sin \varphi F_y
	= \frac{x}{r} F_x + \frac{y}{r} F_y,
	\qquad
	F_{\varphi} = - \sin \varphi F_x + \cos \varphi F_y
	= - \frac{y}{r} F_x + \frac{x}{r} F_y
.\]
\begin{proof}
	תחילה נבחין כי מהגדרתם,
	\[
		x^2 + y^2
		= r^2 \cos^2 \varphi + r^2 \sin^2 \varphi
		= r^2
		\implies \lVert (x, y) \rVert = r
	.\]
	ולכן,
	\[
		F_r
		= \hat{r}(F)
		= \sqrt{F_x^2 + F_y^2}
	.\]
\end{proof}

\end{document}
